\section*{Finite-horizon information, cohort ambiguity, and net inventory-service value}

\paragraph{Status and contribution boundary.}
This appendix supplies a restricted analytic separation and its financial
consequence. The proof below is written in full; exact rational and numerical
checks accompany it. It has not received independent proof review and is not
a Lean formalization. Its priority is unconfirmed. Control without unique
identification, task-directed exploration, Gaussian functional estimation,
moment matching, and the value of information are established ideas
\citep{fhc:waarde,fhc:wagenmaker,fhc:donoho,fhc:heinrich,fhc:blackwell}.
The candidate extension is their explicit combination for a finite-horizon
inventory commitment, with a hard stress budget and the information bill
deducted from both capital and profit. It does not close the broader open
theorem of the preceding drafts.

\subsection*{The observation contract and the financial institution}
Fix $H\ge4$, $0\le a<1$, $\kappa>0$ and
$0\le Q\le\overline Q<(1-a)/\kappa$. There are arbitrarily many finite
cohorts, with $q_i\ge0$, $\sum_iq_i=Q$, and $0<\eta_i\le1$:
\[
 x_{t+1}=a x_t+\kappa\bigl(\sum_iq_i y_{i,t}+u_t\bigr),\qquad
 y_{i,t+1}=(1-\eta_i)y_{i,t}+\eta_i x_{t+1}.
\]
Initial states are reset to zero. Let $h_0,\ldots,h_{H-1}$ be the price
response to a unit impulse, with $h_0=\kappa$, and put
$G=\sum_{t<H}h_t$. Positivity implies $h_t\ge0$.
No positive lower bound on $\eta_i$ is imposed.

The observer purchases $N$ independent reset experiments, each lasting $H$
model periods and returning all $H$ taps with independent Gaussian errors of
known variance $\sigma^2$. Thus the tap means are
$\overline h_t=h_t+\sigma Z_t/\sqrt N$, with independent standard normal
$Z_t$. The same observations are available for identifying the cohorts and
for choosing the financial position. One experiment is not one scalar reading:
it returns $H$ readings and costs $c_H$, including the declared operating and
time charges. Total information cost is $F=Nc_H$. This observation contract
is not an ordinary exchange tape; longer horizons cost both time and readings.
Plant parameters, the feedback cap, reset, and stationarity are assumptions.
Measurements are acquired ahead of a common future service-start deadline.
Cash can earn the common reference rate during acquisition; the model does
not include earlier service opportunities. The stated measurement price is
exogenous, not a calibrated self-impact cost for actual market probes.

The institution is a price-taking, unlevered inventory-warehousing desk.
It buys $I\ge0$ units at $P_0$, holds them for $H$ periods, and closes.
It earns a prequoted service fee $pI$ \emph{after} settlement. Round-trip
execution and carry costs are $bI+dI^2$, where $b\ge0$ and $d>0$.
For example, $b=b_{\rm trade}+P_0 r_f H/(365\cdot24\cdot60)$ for
one-minute model periods. The service premium is unavailable as collateral
until after settlement. Both trading costs and $F$ must be prefunded.
Initial cash is $W$, the stressed-settlement loss budget is $B$, and
$|u_t|\le\overline u$, with $\overline u>0$. Assume
$P_0>\overline u\,\overline G$, so the asset price remains positive.
This desk does not alter market feedback, depth, or another participant's
behaviour. The service fee is exogenous, not an equilibrium price.

For the expected-profit objective, future shocks have conditional mean zero
given the observations. Consequently the conditional expected trading price
change is zero. The financial objective, including every declared charge, is
\[
 V(I,F)=g(I)-F,\qquad g(I)=(p-b)I-dI^2.
\]
This is expected service profit with a hard stress constraint, not a Kelly
growth theorem or a return-prediction result. There are no additional training
or validation tapes, no leverage, uncertain fills, or changing inventory.

\begin{lemma}[Sharp prior bounds and exact prefunded capacity]\label{fhc:capacity}
Write
\[
 \underline G=\kappa\sum_{t<H}a^t,\qquad
 \overline G=\kappa\sum_{t<H}(a+\kappa\overline Q)^t.
\]
Then $\underline G\le G\le\overline G$, and both endpoints are attained
in the declared class. For a gain bound $C\ge G$ and
$0\le F\le\min(B,W)$, define
\[
 \begin{gathered}
 R(v,r)=\frac{2v}{r+\sqrt{r^2+4dv}},\qquad I_u=\frac{(p-b)_+}{2d},\\
 I(C,F)=\min\{I_u,R(B-F,\overline u C+b),R(W-F,P_0+b)\}.
 \end{gathered}
\]
This is the optimal expected-profit position subject to the stated stress
reserve and prefunding constraints. The zero-violation no-observation robust optimum is
$I_{\rm nd}=I(\overline G,0)$, not an arbitrarily conservative benchmark.
The response-informed oracle is $I_\circ=I(G,0)$; it knows no future shocks.
\end{lemma}

\begin{proof}
Let $k_j=\sum_iq_i\eta_i(1-\eta_i)^j$. The normalized price response is
the renewal sequence with defective interarrival probabilities
$s_1=a+\kappa k_0$ and $s_j=\kappa k_{j-1}$ for $j\ge2$.
Their total mass is $a+\kappa Q<1$, with the remaining mass placed at
infinity. The cumulative interarrival distribution at integer $j\ge1$ is
\[
 a+\kappa\sum_iq_i[1-(1-\eta_i)^j].
\]
It lies between $a$ and $a+\kappa\overline Q$. Coupling interarrivals by
the same independent uniforms orders every renewal time and hence the
number of renewals before $H$. Taking expectations proves the two bounds.
The lower endpoint is the no-feedback economy; the upper is one cohort
with $Q=\overline Q$ and $\eta=1$. This also proves sharpness of the
no-observation stress bound.

At date $t\le H$, the price change is
$x_t=\sum_{j<t}h_j u_{t-1-j}$, so
$|x_t|\le\overline u\sum_{j<t}h_j\le\overline uG$.
All costs paid by that date are at most the total prefunded costs.
The constant adverse sequence $u_t=-\overline u$ attains the total price
loss $I\overline uG$ at settlement. Before the service fee is credited,
worst loss from initial wealth is therefore exactly
$I\overline uG+bI+dI^2+F$. The declared robust reserve requirement is
$I\overline uC+bI+dI^2+F\le B$; prefunding requires
$P_0I+bI+dI^2+F\le W$. These increasing quadratics yield the two roots.
Finally $g$ increases up to $I_u$ and decreases after it, proving the
optimizer and the two benchmarks. This exact reserve contract is not a
peak-to-trough drawdown or a repeated-deployment guarantee.
\end{proof}

\begin{lemma}[No data at the same error budget]\label{fhc:nodata-chance}
Let $I_f=\min\{I_u,R(W,P_0+b)\}$ and allow a data-free rule to randomize,
with probability of a stress-budget violation at most $\gamma$ in every
economy of the class. Its greatest expected service profit is
\[
 V_{{\rm nd},\gamma}=(1-\gamma)g(I_{\rm nd})+\gamma g(I_f).
\]
It is attained by the displayed two-position lottery. This is a theoretical
benchmark, not a recommended trading policy. The corresponding response-
informed same-error-budget oracle has value
$(1-\gamma)g(I_\circ)+\gamma g(I_f)$.
\end{lemma}

\begin{proof}
Truncating an action at $I_f$ weakly improves profit and lowers risk, so
restrict to $0\le I\le I_f$. At the attainable worst economy with gain
$\overline G$, every $I>I_{\rm nd}$ is unsafe unless $I_{\rm nd}=I_f$,
in which case all larger actions are already dominated. Thus a feasible
lottery puts at most $\gamma$ mass above $I_{\rm nd}$. Below that point
profit is at most $g(I_{\rm nd})$, and above it at most $g(I_f)$.
This gives the upper bound. The two-point lottery attains it and is unsafe
in any economy with probability at most $\gamma$. Replace $\overline G$
by the known $G$ for the oracle statement.
\end{proof}

\begin{theorem}[Uniform cost-aware finite-horizon decision]\label{fhc:decision}
Fix $0<\gamma<1/2$, $z=\Phi^{-1}(1-\gamma)$,
$s_N=\sigma\sqrt{H/N}$, and
\[
 Y=\sum_{t<H}\overline h_t,\qquad
 C_N=\min\{\overline G,\max\{\underline G,Y+zs_N\}\}.
\]
Purchase a fixed, affordable $N$ and take $I_N=I(C_N,F)$.
The following hold uniformly over the declared positive-cohort class,
without dependence on cohort spacing, cohort count, or a memory-rate floor.
\begin{enumerate}
\item With probability at least $1-\gamma$ over measurements, the position
      satisfies the stress budget for \emph{every} bounded future shock path.
\item Put $r_0=\overline u\underline G+b$, $I_{\max}=W/P_0$,
\[
 \begin{aligned}
 L_G&=\frac{\overline u I_{\max}}{r_0}, &
 L_F&=\frac{1}{\min(r_0,P_0+b)},\\
 A&=2z\sigma\sqrt H\,(p-b)_+L_G, & D&=1+(p-b)_+L_F.
 \end{aligned}
\]
With probability at least $1-2\gamma$,
\[
 0\le g(I_\circ)-V(I_N,F)\le A N^{-1/2}+D c_HN.
\]
For sufficiently small $c_H>0$, choosing
$N=\lceil(A/(2Dc_H))^{2/3}\rceil$ is affordable and gives a bound
$3\,2^{-2/3}A^{2/3}(Dc_H)^{1/3}+Dc_H$, when $A>0$.
This is an upper bound from this procedure, not a new minimax optimality claim.
\item Let $I_-=I(G+2zs_N,F)$. A sufficient condition for strictly positive
      expected net information value relative to the strongest no-data rule
      at the \emph{same} error budget is
\[
 \underbrace{(1-2\gamma)g(I_-)-F-V_{{\rm nd},\gamma}}_{\displaystyle\mathcal L_N}>0.
\]
Precisely, $\mathbb E_D[V(I_N,F)]-V_{{\rm nd},\gamma}\ge\mathcal L_N$.
The conclusion can fail at high data costs or a low service premium.
\end{enumerate}
\end{theorem}

\begin{proof}
$Y-G$ is Gaussian with variance $H\sigma^2/N$. On
$\{Y-G\ge-zs_N\}$, which has probability $1-\gamma$, $C_N\ge G$;
clipping preserves this because $G$ lies between the prior bounds.
Lemma~\ref{fhc:capacity} proves simultaneous safety for every future
bounded path on this single event, including the data-selected position.
There is no tail beyond $H$ to estimate, explaining the absence of a
memory-floor assumption. This does not certify an infinite horizon.

On $E=\{|Y-G|\le zs_N\}$, of probability $1-2\gamma$,
$G\le C_N\le G+2zs_N$. The position function is decreasing in each
of $C$ and $F$, so $I_-\le I_N\le I_\circ$ on $E$.
To bound the position difference, implicitly differentiate the stress root:
its derivatives are $-\overline u I/(\overline u C+b+2dI)$ in $C$
and $-1/(\overline u C+b+2dI)$ in $F$.
Clip this root at $I_{\max}$ without changing the final minimum, since
the prefunding root is at most $I_{\max}$. The clipped root is
$L_G$-Lipschitz in $C$ and $1/r_0$-Lipschitz in $F$.
The prefunding root is independent of $C$ and
$1/(P_0+b)$-Lipschitz in $F$. Minimum with $I_u$ preserves these
common Lipschitz bounds. Hence
\[
 0\le I_\circ-I_N\le 2L_Gzs_N+L_FF.
\]
On $[0,I_u]$, $0\le g'\le(p-b)_+$. Multiplication and addition of
the paid fee give the regret bound. Minimizing
$AN^{-1/2}+Dc_HN$ over positive real $N$ gives the stated root and
value. Rounding up adds at most $Dc_H$; affordability follows as
$c_HN=O(c_H^{1/3})\to0$. The case $A=0$ has no positive-gain
information task and does not need this optimization.

On $E$, $g(I_N)\ge g(I_-)$, while $g(I_N)\ge0$ always because
$0\le I_N\le I_u$. Thus
$\mathbb E_D g(I_N)\ge(1-2\gamma)g(I_-)$.
Subtracting $F$ and the same-error-budget no-data value proves the last assertion.
Conditional mean-zero future flows justify the stated expected-profit
interpretation; their distribution is not learned from extra tapes.
\end{proof}

\begin{theorem}[A two-cohort identification obstruction under the same observations]\label{fhc:pair}
Fix $Q>0$, $a+\kappa Q<1$, and $0<\eta_0<1$. For sufficiently small
$\delta>0$, compare the two economies
\[
\begin{array}{c|cc}
 &\text{normalized weights}&\text{adaptation rates}\\ \hline
\mathrm A_\delta &(1/2,1/2)&(\eta_0-\delta,\eta_0+\delta)\\
\mathrm B_\delta &(1/5,4/5)&(\eta_0-2\delta,\eta_0+\delta/2).
\end{array}
\]
Actual weights are $Q$ times those displayed. For fixed $H\ge4$,
$D_H(\delta)=\|h^{\rm B}_{0:H-1}-h^{\rm A}_{0:H-1}\|_2
=\Theta(\delta^3)$. The exact minimax error for distinguishing these
two \emph{known} alternatives from $N$ reset experiments is
\[
 e_N=\Phi\!\left(-\frac{\sqrt N D_H(\delta)}{2\sigma}\right).
\]
Therefore error at most $\beta<1/2$ requires
$N\ge[2\sigma\Phi^{-1}(1-\beta)/D_H(\delta)]^2
=\Theta(\delta^{-6})$.
In particular any recovery of the two normalized weights to sup-norm error
less than $0.15$ with probability at least $1-\beta$ in both worlds must
satisfy this testing requirement.
Both worlds contain exactly two distinct rates. Both have identical infinite
summed response, but their finite-horizon gains are not asserted identical.
\end{theorem}

\begin{proof}
The two rate distributions have equal mass, mean $\eta_0$, and variance
$\delta^2$. Their third centered moments are respectively $0$ and
$-3\delta^3/2$. Thus their kernels $k_j$ agree at $j=0,1$ and
$k^{\rm B}_2-k^{\rm A}_2=-3Q\delta^3/2$. The response recursion
gives equal $h_0,h_1,h_2$ and
\[
 h^{\rm B}_3-h^{\rm A}_3=-3\kappa^2Q\delta^3/2,
\]
which is a lower bound for $D_H$ in absolute value.

For an upper bound put $f_t(\eta)=\eta(1-\eta)^t$.
Its third derivative is
$3t(t-1)(1-\eta)^{t-2}
-\eta t(t-1)(t-2)(1-\eta)^{t-3}$, with absent low-degree terms zero.
The sum over $t\ge0$ of its absolute value is at most $12/\eta^3$.
For $\delta\le\eta_0/4$, all intervening rates are at least
$v=\eta_0/2$. Taylor's integral remainder in $\ell_1$ is therefore at
most $2|\eta-\eta_0|^3/v^3$. The matching first two moments cancel
the quadratic Taylor terms. The sum of the two third absolute centered
moments is $(1+17/10)\delta^3=27\delta^3/10$, giving
\[
 \|k^{\rm B}-k^{\rm A}\|_1\le\frac{27Q}{5v^3}\delta^3.
\]
If $K(z)=\sum k_tz^t$ and $H(z)=\sum h_tz^t$, the recursion gives
$H(z)=\kappa/[1-az-\kappa zK(z)]$. Consequently
$H_{\rm B}-H_{\rm A}=zH_{\rm A}H_{\rm B}(K_{\rm B}-K_{\rm A})$.
The infinite response sums of both economies equal
$S=\kappa/(1-a-\kappa Q)$. Young's inequality yields
$D_H\le S^2\|k^{\rm B}-k^{\rm A}\|_1=O(\delta^3)$.
This proves the two-sided order with constants independent of $\delta$.

For equal-covariance Gaussian laws, the likelihood-ratio test projects the
sample mean onto the difference of the two known means, and thresholds at
their midpoint. The projected normal law gives the displayed error in each
world. Its equal-prior Bayes optimality and symmetry establish minimax
optimality. A weight estimator accurate within $0.15$ distinguishes the
two disjoint normalized-weight balls, proving the identification lower bound.
Finite-horizon gains need not be equal: only their difference is of cubic
order. The test concerns latent cohort recovery, not identification of the
input-output response vector, and does not cover alternative probe designs.
\end{proof}

\begin{corollary}[A priced separation]\label{fhc:priced}
Fix $c_H>0$ and an affordable $N$ for which $\mathcal L_N>0$ at the limiting merged
economy $\delta=0$. For all sufficiently small positive $\delta$, the
financial rule of Theorem~\ref{fhc:decision} has a
strictly positive expected net information value, with its stated safety
probability. Its observation requirement has no dependence on $\delta$.
At this same fixed information price, sufficiently small $\delta$ makes cohort
identification at error $\beta$ unaffordable within the same hard fee budget
$B$: $c_H[2\sigma\Phi^{-1}(1-\beta)/D_H(\delta)]^2>B$.
Any procedure limited to that many affordable observations is covered,
including one that stops early, since it can be viewed as a function of the
full affordable sample. A rational agent may abandon identification and use
the no-data rule. There is no assertion that regulation requires identification.
\end{corollary}

\begin{proof}
Apply the two theorems. The first statement uses the strictly positive
lower bound on net value, rather than an assumption that more information
is automatically beneficial. The second follows from $D_H=\Theta(\delta^3)$.
For a neighbourhood of a fixed $\delta$ where the strict value inequality
holds, continuity preserves it. To claim the inequality for all sufficiently
small $\delta$, it suffices that it hold strictly at the limiting merged
economy: finite $G_H$ and $I_-$ are continuous there. The numerical example
reports both the primary pair and that limiting check. No finite-class grid
is used in the uniform decision guarantee.
\end{proof}

\paragraph{What the result does not establish.}
The direct gain-sum certificate is exactly ordinary Gaussian projection onto
the vector of ones. It is not a new estimator and does not beat that benchmark.
The latent-cohort lower bound is not a lower bound for learning the observable
response vector; comparing it with FIR model-identification results for
$\mathcal H_\infty$ norms would conflate different tasks \citep{fhc:tu}.
No optimal estimation rate over all observation designs is claimed.
The data cost is not an empirically measured exchange probing cost.
An ordinary public-tape bridge, endogenous quotes and depth, incomplete resets,
unknown or misspecified noise, and independent replication remain open.
The information purchase in the example is predeclared. Its positive value
is pointwise at the displayed economy; no prior over unknown economies or
globally optimal Bayesian/robust acquisition policy is claimed.
The fixed horizon, independent measurement contract, and conditional mean-zero
future flows cannot be removed by the simulation. At another horizon the
certificate must be recomputed. A rolling sequence needs simultaneous or
time-uniform coverage; a one-decision guarantee is insufficient.

\paragraph{Literature priority.}
The reviewed sources already establish the general principles used here.
The explicit dollar-cost/capital formulation and the combined result above
are candidate additions to this project, not verified first-in-literature
claims. A specialist must check whether they are an immediate corollary of
existing task-specific or constrained-decision results before publication.

\begin{thebibliography}{9}\small
\bibitem[van Waarde et~al.(2020)]{fhc:waarde} H.~J. van Waarde, J. Eising,
H.~L. Trentelman and M.~K. Camlibel. Data Informativity: A New Perspective on
Data-Driven Analysis and Control. \emph{IEEE TAC} 65, 4753--4768.
\bibitem[Wagenmaker et~al.(2021)]{fhc:wagenmaker} A.~J. Wagenmaker,
M. Simchowitz and K. Jamieson. Task-Optimal Exploration in Linear Dynamical
Systems. \emph{ICML}, PMLR 139, 10641--10652.
\bibitem[Donoho(1994)]{fhc:donoho} D.~L. Donoho. Statistical Estimation and
Optimal Recovery. \emph{Annals of Statistics} 22, 238--270.
\bibitem[Heinrich and Kahn(2018)]{fhc:heinrich} P. Heinrich and J. Kahn.
Strong Identifiability and Optimal Minimax Rates for Finite Mixture Estimation.
\emph{Annals of Statistics} 46, 2844--2870.
\bibitem[Blackwell(1953)]{fhc:blackwell} D. Blackwell. Equivalent Comparisons
of Experiments. \emph{Annals of Mathematical Statistics} 24, 265--272.
\bibitem[Tu et~al.(2018)]{fhc:tu} S. Tu, R. Boczar and B. Recht.
Minimax Lower Bounds for $\mathcal H_\infty$-Norm Estimation.
\emph{arXiv:1809.10855}.
\end{thebibliography}
